% chapters/appendix-ch05-sql.tex
\chapter{SQL Appendix: Chapter 5 (Normalization)}

\section{DDL (normalized)}
\begin{verbatim}
CREATE TABLE doctor (
  doctor_id INT PRIMARY KEY AUTO_INCREMENT,
  name VARCHAR(120) NOT NULL,
  specialty VARCHAR(120) NOT NULL
) ENGINE=InnoDB;

CREATE TABLE patient (
  patient_id INT PRIMARY KEY AUTO_INCREMENT,
  name VARCHAR(120) NOT NULL,
  birth_date DATE NOT NULL
) ENGINE=InnoDB;

CREATE TABLE appointment (
  appointment_id INT PRIMARY KEY AUTO_INCREMENT,
  start_time DATETIME NOT NULL,
  doctor_id INT NOT NULL,
  patient_id INT NOT NULL,
  CONSTRAINT fk_appt_doctor FOREIGN KEY (doctor_id) REFERENCES doctor(doctor_id),
  CONSTRAINT fk_appt_patient FOREIGN KEY (patient_id) REFERENCES patient(patient_id)
) ENGINE=InnoDB;
\end{verbatim}

\section{Sample data}
\begin{verbatim}
INSERT INTO doctor (name, specialty) VALUES
  ('Joana Reis', 'cardiology'),
  ('Hugo Pereira', 'pediatrics');

INSERT INTO patient (name, birth_date) VALUES
  ('Ana Costa', '1990-04-11'),
  ('Rui Silva', '1985-09-03');

INSERT INTO appointment (start_time, doctor_id, patient_id) VALUES
  ('2026-02-03 09:00:00', 1, 1),
  ('2026-02-03 09:30:00', 1, 2);
\end{verbatim}

\section{Example queries}
\begin{verbatim}
-- Q1: Appointments with doctor and patient names
SELECT a.appointment_id, a.start_time, d.name AS doctor, p.name AS patient
FROM appointment a
JOIN doctor d ON d.doctor_id = a.doctor_id
JOIN patient p ON p.patient_id = a.patient_id;

-- Q2: Detect redundant specialty storage (should be in doctor only)
SELECT a.appointment_id, d.specialty
FROM appointment a
JOIN doctor d ON d.doctor_id = a.doctor_id;
\end{verbatim}
